两块无限宽静止平板位于 \(y=\pm h\)，中间充满密度 \(\rho>0\)、动力黏度 \(\mu>0\) 的牛顿流体。定义运动黏度 \(\nu=\mu/\rho\)，假定速度场 \(\boldsymbol u=u(y,t)\boldsymbol e_x\) 且沿流向均匀。初态 \(u(y,0)=0\)，从 \(t=0\) 起沿 \(x\) 方向施加恒定每单位质量体力 \(F>0\)（等效于恒定压强梯度）；两壁严格无滑移。总流量 \(Q(t)=\int_{-h}^{h}u(y,t)\,dy\) 取单位跨距。
(1) 从完整Navier--Stokes方程严格约化出一维初边值问题；给出稳态剖面、瞬态的全部偶对称本征模和启动过程的精确余弦级数。明确其收敛与初始条件。
(2) 由正交性求精确的总流量级数 \(Q(t)\)、稳态流量和最慢衰减时间；证明 \(Q(0)=0\)、\(Q(t)<Q_\infty\) 且流量从静止单调增加，不能仅用稳态近似代替瞬态。
(3) 计算两壁剪应力的时间演化；从控制方程直接推导动能变化、外力功和黏滞耗散的逐时能量恒等式，核对长期稳态功率收支与 \(t=0^+\) 的零剪应力极限。
(4) 在已达到恒定稳态以后叠加小幅谐波加速度 \(F_1\operatorname{Re}e^{i\Omega t}\)，\(\Omega>0\)。求受迫速度复振幅、流量复导纳及低频/高频极限；证明平均外力输入功率等于平均黏滞热耗散，并解释频率相关的壁面边界层厚度。